\chapter{Supporting Checks}

This appendix contains the checks most relevant to interpreting the main
results. Additional checks are documented separately.

\section{Accident Characteristics}
\label{app:accident-characteristics}

Figures~\ref{fig:vehicles-per-accident} and
\ref{fig:accident-types-severity} complement the hour, season and daylight
counts in the data chapter. They describe the selected injury accidents without
a traffic-exposure denominator.

\begin{figure}[htbp]
\centering
\includegraphics[width=0.72\textwidth]{vehicles_per_accident.pdf}
\caption[Registered vehicles involved in rural injury accidents.]{Number of registered vehicles involved in the 6,414 rural injury
accidents. Counts describe the selected sample, not risk by vehicle count.}
\label{fig:vehicles-per-accident}
\end{figure}

\begin{figure}[htbp]
\centering
\includegraphics[width=\textwidth]{accident_types_by_severity.pdf}
\caption[Accident-type composition by injury severity.]{Accident-type composition for minor injury (blue, code 3) and
serious or fatal injury (red, codes 1--2). Percentages sum to 100 within
each separate severity group.}
\label{fig:accident-types-severity}
\end{figure}

\clearpage
\section{Supporting Statistical Models}
\label{app:statistical-models}

This section gives the fitting details for the supporting matched-time and severity models.

The matched-time models use \texttt{statsmodels} version 0.14.6 and
\texttt{ConditionalLogit}, maximising the conditional logistic likelihood by
BFGS optimisation. Conditioning on the single case in each stratum eliminates
its intercept. Model-based covariance is the inverse negative Hessian of the
conditional log-likelihood; ORs and normal-Wald intervals follow
Equation~\ref{eq:odds-ratio-ci}.

The joint model includes mean-wind and temperature categories together, using
these same references and matched times with both measures available. The
daylight model compares darkness and civil twilight with daylight at the
accident location, under the same calendar matching and fitting procedure.
The seasonal model uses mean-wind groups 0--10 (reference), 10--15, and at least
15 m/s, with Winter as the reference season for wind interactions. A
likelihood-ratio test compares models with and without the six interaction
terms against a chi-squared distribution with six degrees of freedom.
Season-specific log-odds contrasts combine the relevant coefficients; their
standard errors use the full covariance matrix, with intervals computed as in
Equation~\ref{eq:odds-ratio-ci} using 1.96.

The adjusted severity model uses \texttt{statsmodels.Logit}, estimating an
intercept and categorical coefficients jointly by Bernoulli maximum
likelihood with Newton optimisation. Among injury accidents with qualifying
wind, temperature, and daylight, the outcome is one for serious or fatal
injury (codes 1--2) and zero for minor injury (code 3). Reference categories
are mean wind 0--10 m/s, temperature 0--3\(^{\circ}\)C, daylight, 10:00--16:00,
and Summer. The fit requests HC1 heteroskedasticity-robust covariance;
adjusted ORs and normal-Wald intervals follow
Equation~\ref{eq:odds-ratio-ci} with those robust standard errors.
These ORs concern severity conditional on an injury accident, not accident
occurrence.

Table~\ref{tab:severity-conditions} shows the adjusted comparison of
serious or fatal and minor-injury accidents. It describes severity among
recorded injury accidents; it does not estimate accident occurrence or risk
per vehicle-kilometre.

\begin{table}[H]
\centering
\footnotesize
\caption{Adjusted odds ratios for a serious-or-fatal outcome among recorded injury accidents. All listed variables are included in the same model.}
\label{tab:severity-conditions}
\begin{tabularx}{\textwidth}{L{0.20\textwidth}L{0.45\textwidth}X}
\toprule
Variable & Comparison & Adjusted OR (95\% CI) \\
\midrule
Mean wind & 10-15 m/s vs 0-10 m/s & 0.99 (0.82--1.21) \\ \grayhline
Mean wind & $\geq$15 m/s vs 0-10 m/s & 0.93 (0.69--1.26) \\ \grayhline
Temperature & <-6 °C vs 0-3 °C & 1.13 (0.79--1.62) \\ \grayhline
Temperature & -6--3 °C vs 0-3 °C & 0.85 (0.63--1.15) \\ \grayhline
Temperature & -3-0 °C vs 0-3 °C & 0.83 (0.66--1.04) \\ \grayhline
Temperature & 3-6 °C vs 0-3 °C & 1.11 (0.88--1.39) \\ \grayhline
Temperature & 6-9 °C vs 0-3 °C & 1.41 (1.10--1.79) \\ \grayhline
Temperature & 9-12 °C vs 0-3 °C & 1.28 (0.98--1.68) \\ \grayhline
Temperature & 12-15 °C vs 0-3 °C & 1.33 (0.98--1.81) \\ \grayhline
Temperature & $\geq$15 °C vs 0-3 °C & 1.21 (0.83--1.76) \\ \grayhline
Daylight & Darkness vs Daylight & 0.89 (0.71--1.12) \\ \grayhline
Daylight & Civil twilight vs Daylight & 1.09 (0.87--1.36) \\ \grayhline
Time of day & 00-06 vs 10-16 & 0.86 (0.64--1.16) \\ \grayhline
Time of day & 06-10 vs 10-16 & 0.65 (0.52--0.82) \\ \grayhline
Time of day & 16-20 vs 10-16 & 0.98 (0.84--1.14) \\ \grayhline
Time of day & 20-24 vs 10-16 & 0.77 (0.61--0.98) \\ \grayhline
Season & Winter vs Summer & 0.82 (0.64--1.05) \\ \grayhline
Season & Spring vs Summer & 0.96 (0.78--1.18) \\ \grayhline
Season & Autumn vs Summer & 0.83 (0.65--1.05) \\
\bottomrule
\end{tabularx}
\end{table}


\clearpage
\section{Period and Counter-Era Sensitivity}
\label{app:period-sensitivity}

These checks assess whether period choice changes the upper-wind pattern.

Figures~\ref{fig:oe-early-period} and \ref{fig:oe-counter-period} repeat the
weather-frequency comparison for 2007--2018 and 2019--2024. They are supporting
All year checks; upper-bin counts should be read alongside the ratios.

\begin{figure}[!htbp]
\centering
\includegraphics[width=\textwidth,height=0.65\textheight,keepaspectratio]{weather_oe_2007_2018.pdf}
\caption[All year weather-frequency O/E, 2007--2018.]{All year weather-frequency O/E for rural injury accidents, 2007--2018, with
accidents and background weather restricted to the same period. Blue and red
show minor and serious or fatal injury; labels give counts and the line marks
O/E = 1.}
\label{fig:oe-early-period}
\end{figure}

\begin{figure}[p]
\centering
\includegraphics[width=\textwidth,height=0.73\textheight,keepaspectratio]{weather_oe_2019_2024.pdf}
\caption[All year weather-frequency O/E, 2019--2024.]{All year weather-frequency O/E for rural injury accidents, 2019--2024, with
accidents and background weather restricted to the same period. Blue and red
show minor and serious or fatal injury; labels give counts and the line marks
O/E = 1.}
\label{fig:oe-counter-period}
\end{figure}

Figure~\ref{fig:corrected-counter-period} applies the traffic correction within
2019--2024, allowing comparison with the uncorrected counter-era population in
Figure~\ref{fig:oe-counter-period}. Its temperature panel is corrected
weather-frequency O/E, not a VKT rate.

\begin{figure}[p]
\centering
\includegraphics[width=\textwidth,height=0.73\textheight,keepaspectratio]{weather_oe_traffic_corrected_2019_2024.pdf}
\caption[Traffic-corrected All year O/E by injury group, 2019--2024.]{Approximately traffic-corrected O/E for rural injury accidents, 2019--2024,
All year. Blue and red show minor and serious or fatal injury using reweighted
expected counts; labels give observed accidents and the line marks O/E = 1.}
\label{fig:corrected-counter-period}
\end{figure}
\clearpage

\section{Supporting Allocated Daily-Counter Analysis}
\label{app:allocated-daily-counter}

This analysis checks the wind association under an alternative counter-linked
exposure construction. The supporting model uses observed 24-hour traffic
totals and allocates them across wind intervals according to each day's valid
10-minute weather observations. Directional channels are first summed to one
count for each physical counter and date. Accidents require a counter on the
exact registered road section and year, positive traffic on the accident date,
and a qualifying accident-time mean-wind observation. Eligible counter-days
require at least 108 valid 10-minute observations over the full day. The
allocated traffic is

\begin{equation}
\label{eq:allocated-counter-traffic}
\widehat Q_{cdj}=Q_{cd}\frac{N_{cdj}}{N_{cd}},
\end{equation}

where \(c\) indexes physical counters, \(d\) dates, and \(j\) wind intervals;
\(Q_{cd}\) is observed 24-hour traffic, \(N_{cdj}\) counts valid 10-minute
observations in interval \(j\), \(N_{cd}\) counts all valid observations that
day, and \(\widehat Q_{cdj}\) is traffic allocated to that interval.
Accidents use event-time mean wind from the denominator's station.
The allocation assumes uniform traffic across valid observation times.

Traffic and accidents are aggregated by counter, year, and wind interval.
The model uses \texttt{ConditionalPoisson} from \texttt{statsmodels} to estimate
coefficients for
10--15 and at least 15 m/s relative to 0--10 m/s. BFGS maximises the Poisson
likelihood conditional on each counter--year's total accident count, removing
its intercept. The offset is the logarithm of summed allocated traffic. An offset has its
coefficient fixed at one: it accounts for differing exposure so coefficients
compare accident rates rather than counts. Here exposure is allocated vehicle
counts, not vehicle-kilometres. Rate ratios are \(\exp(\widehat\beta)\), where
\(\widehat\beta\) is the fitted log-rate contrast. Normal-Wald 95\% intervals
use the same exponentiated-coefficient construction as
Equation~\ref{eq:odds-ratio-ci}, with standard errors from the inverse negative
Hessian of the conditional log-likelihood.

This supporting model is separate from the 694-accident main VKT analysis. The
tables below show how the 2019--2024 rural injury sample narrows under its
linkage. The 767 records at the last broad selection step include five records
that do not enter the final allocated-rate model, which retains 762 accidents.

Table~\ref{tab:allocated-selection} traces the linkage exclusions for this
supporting sample.

\begin{table}[H]
\centering
\small
\caption{Selection of accidents for the observed daily-counter analysis}

\begin{tabularx}{0.86\textwidth}{Xrr}
\toprule
Selection step & Accidents & Share of 2019--2024 sample \\
\midrule
Rural injury accidents in 2019--2024 & 1,863 & 100.0\% \\ \grayhline
Exact road-section/year counter candidate & 864 & 46.4\% \\ \grayhline
Assigned counter within 20 km & 863 & 46.3\% \\ \grayhline
Positive count and valid full-day wind & 767 & 41.2\% \\
\bottomrule
\end{tabularx}
\end{table}


Table~\ref{tab:allocated-sample} compares retained accidents with the broader
sample and excluded groups by severity, vehicles and season, making selection
differences visible.

\begin{table}[H]
\centering
\footnotesize
\caption{Characteristics of accidents retained and excluded by the daily-counter linkage}

\begin{tabularx}{\textwidth}{Xrrrr}
\toprule
Characteristic & All & No link & Linked, excluded & Retained \\
\midrule
Accidents & 1,863 & 999 & 106 & 758 \\ \grayhline
Share of all accidents & 100.0\% & 53.6\% & 5.7\% & 40.7\% \\ \grayhline
Distinct road sections & 563 & 374 & 59 & 189 \\ \grayhline
Serious or fatal & 25.6\% & 29.4\% & 17.9\% & 21.6\% \\ \grayhline
One vehicle & 67.7\% & 75.2\% & 54.7\% & 59.8\% \\ \grayhline
Winter & 32.6\% & 28.0\% & 41.5\% & 37.3\% \\ \grayhline
Spring & 13.0\% & 11.7\% & 8.5\% & 15.4\% \\ \grayhline
Summer & 36.0\% & 41.8\% & 33.0\% & 28.6\% \\ \grayhline
Autumn & 18.4\% & 18.4\% & 17.0\% & 18.6\% \\
\bottomrule
\end{tabularx}
\end{table}


The geographic selectivity of this traffic-linked sample is documented in
Figure~\ref{fig:counter-coverage}.

\begin{figure}[H]
\centering
\includegraphics[width=0.9\textwidth]{counter_coverage.pdf}
\caption[Coverage of the supporting allocated daily-counter sample.]{Geographic coverage of located daily counters and rural injury
accidents retained or excluded by the daily-counter linkage, 2019--2024.}
\label{fig:counter-coverage}
\end{figure}

\section{Traffic Sensitivity Checks}
\label{app:traffic-sensitivity}

These checks assess sensitivity to counter assignment and recorded zero traffic.

Table~\ref{tab:allocated-radius} is a separate daily-mean-wind check under
5, 10, and 20 km accident-to-counter assignment limits. It classifies each
whole day's accident count and observed traffic by that day's mean wind,
without allocation among 10-minute weather intervals. The conditional Poisson
fit otherwise uses the counter--year strata, 0--10 m/s reference, BFGS
estimation, and RR/CI calculation described above; its offset is log summed
observed daily traffic in each counter--year--wind group. The upper bin is at
least 15 m/s. Included accidents count the linked sample before model filtering.

\begin{table}[H]
\centering
\small
\caption[Assignment-distance sensitivity of the daily-mean-wind model.]{Upper-bin estimate from the daily-mean-wind counter model under three accident-to-counter assignment distances.}
\label{tab:allocated-radius}

\begin{tabular}{rrrrr}
\toprule
Maximum distance & Included accidents & Upper-bin accidents & Rate ratio & 95\% CI \\
\midrule
5 km & 585 & 12 & 2.91 & 1.58--5.37 \\ \grayhline
10 km & 750 & 14 & 2.78 & 1.58--4.88 \\ \grayhline
20 km & 767 & 14 & 2.73 & 1.55--4.78 \\
\bottomrule
\end{tabular}
\end{table}


Table~\ref{tab:traffic-sensitivity} checks the effect of excluding recorded
zero counter-days from the daily traffic-response estimate. This supporting
check groups whole days by daytime mean wind (10:00--21:59) and compares
summed observed traffic with the counter--year--month--weekday mean expectation,
recomputed for each zero-count policy. Its 95\% intervals are the 2.5th and
97.5th percentiles of 1,000 bootstrap replicates that resample physical counters
with replacement, keeping their observed and expected traffic totals together.
The precomputed calendar expectations remain fixed within each bootstrap.

\begin{table}[H]
\centering
\footnotesize
\caption{Effect of excluding zero counter-days on the daily traffic response.}
\label{tab:traffic-sensitivity}
\begin{tabularx}{\textwidth}{L{0.19\textwidth}L{0.27\textwidth}L{0.22\textwidth}X}
\toprule
Check & Data included & Estimate (95\% interval) & Records \\
\midrule
Daily traffic, 20--25 m/s & Retain zero counter-days & 68.94\% (62.33--74.09) & 1,175 days \\ \grayhline
Daily traffic, 20--25 m/s & Exclude zero counter-days & 69.29\% (62.70--74.52) & 1,090 days \\
\bottomrule
\end{tabularx}
\end{table}

